\documentclass[10pt,a4paper]{amsart}
\usepackage{amsmath,amssymb,amsthm,graphicx}
% Retained preamble for the Stacks Project chapter snapshot in this directory.
% The chapter's own preamble.tex is a copy of the shared Stacks Project preamble; the block
% below is sliced MECHANICALLY, line for line, from it. Nothing was added, removed or
% redefined; the wrapper class, the xr-hyper/hyperref wiring and the externaldocument list
% are omitted because the assembler supplies the document class.
%
% lines taken from preamble.tex: 3, 5, 9, 11, 12, 156, 157, 158, ... (51 lines total)
%
% package coverage restored from the chapter's own preamble (2026-10-13): xy, xy-2cell

\usepackage[all]{xy}
\xyoption{2cell}
\UseAllTwocells

\newenvironment{reference}{\comment}{\endcomment}
\newenvironment{slogan}{\comment}{\endcomment}
\newenvironment{history}{\comment}{\endcomment}
\theoremstyle{plain}
\newtheorem{theorem}[subsection]{Theorem}
\newtheorem{proposition}[subsection]{Proposition}
\newtheorem{lemma}[subsection]{Lemma}
\theoremstyle{definition}
\newtheorem{definition}[subsection]{Definition}
\newtheorem{example}[subsection]{Example}
\newtheorem{exercise}[subsection]{Exercise}
\newtheorem{situation}[subsection]{Situation}
\theoremstyle{remark}
\newtheorem{remark}[subsection]{Remark}
\newtheorem{remarks}[subsection]{Remarks}
\numberwithin{equation}{subsection}
\def\lim{\mathop{\mathrm{lim}}\nolimits}
\def\colim{\mathop{\mathrm{colim}}\nolimits}
\def\Spec{\mathop{\mathrm{Spec}}}
\def\Hom{\mathop{\mathrm{Hom}}\nolimits}
\def\Ext{\mathop{\mathrm{Ext}}\nolimits}
\def\SheafHom{\mathop{\mathcal{H}\!\mathit{om}}\nolimits}
\def\SheafExt{\mathop{\mathcal{E}\!\mathit{xt}}\nolimits}
\def\Sch{\mathit{Sch}}
\def\Mor{\mathop{\mathrm{Mor}}\nolimits}
\def\Ob{\mathop{\mathrm{Ob}}\nolimits}
\def\Sh{\mathop{\mathit{Sh}}\nolimits}
\def\NL{\mathop{N\!L}\nolimits}
\def\CH{\mathop{\mathrm{CH}}\nolimits}
\def\proetale{{pro\text{-}\acute{e}tale}}
\def\etale{{\acute{e}tale}}
\def\QCoh{\mathit{QCoh}}
\def\Ker{\mathop{\mathrm{Ker}}}
\def\Im{\mathop{\mathrm{Im}}}
\def\Coker{\mathop{\mathrm{Coker}}}
\def\Coim{\mathop{\mathrm{Coim}}}
\DeclareMathSymbol{\boxtimes}{\mathbin}{AMSa}{"02}
\def\QCohstack{\mathcal{QC}\!\mathit{oh}}
\def\Cohstack{\mathcal{C}\!\mathit{oh}}
\def\Spacesstack{\mathcal{S}\!\mathit{paces}}
\def\Quotfunctor{\mathrm{Quot}}
\def\Hilbfunctor{\mathrm{Hilb}}
\def\Curvesstack{\mathcal{C}\!\mathit{urves}}
\def\Polarizedstack{\mathcal{P}\!\mathit{olarized}}
\def\Complexesstack{\mathcal{C}\!\mathit{omplexes}}
\def\Pic{\mathop{\mathrm{Pic}}\nolimits}
\def\Picardstack{\mathcal{P}\!\mathit{ic}}
\def\Picardfunctor{\mathrm{Pic}}
\def\Deformationcategory{\mathcal{D}\!\mathit{ef}}
\usepackage{cleveref}
\crefname{originallemma}{original lemma}{original lemmas}
\Crefname{originallemma}{Original lemma}{Original lemmas}
\makeatletter
\newlabel{lemma-space-over-ring}{{08VG}{source}{}{sourcepointer-1}{}}
\newlabel{lemma-space-over-ring@cref}{{[originallemma][0][]08VG}{[1][0][]source}{}{}{}}
\makeatother
\crefname{originalremark}{original remark}{original remarks}
\Crefname{originalremark}{Original remark}{Original remarks}
\makeatletter
\newlabel{remark-compute-L-pi-shriek-spaces}{{08VL}{source}{}{sourcepointer-2}{}}
\newlabel{remark-compute-L-pi-shriek-spaces@cref}{{[originalremark][0][]08VL}{[1][0][]source}{}{}{}}
\makeatother
\crefname{originalremark}{original remark}{original remarks}
\Crefname{originalremark}{Original remark}{Original remarks}
\makeatletter
\newlabel{remark-map-sections-over-U}{{08SX}{source}{}{sourcepointer-3}{}}
\newlabel{remark-map-sections-over-U@cref}{{[originalremark][0][]08SX}{[1][0][]source}{}{}{}}
\makeatother
\title{The cotangent complex of a morphism to a constant algebraic space as a lower shriek of differentials}
\author{The Stacks Project Authors}
\date{stacks-project:cotangent.tex:local-snapshot-20261012; GNU FDL 1.2 or later}
\begin{document}
\maketitle
\setcounter{subsection}{1}
\noindent\textit{Editorial note.} All delivered mathematical text is the author's own, copied from the pinned Stacks Project chapter file cotangent.tex; no statement or proof step was rewritten, reordered or removed.\par\medskip
\noindent\textit{Editorial note.} The statement displays a chain of three equalities; the proof justifies the second and third first and then identifies the first with the cotangent complex of the ring map.\par\medskip
\noindent\textit{Editorial note.} The three references of the counted spans are recorded as non-load-bearing navigation pointers; the argument that consumes them stays inside the counted proof span.\par\medskip
\noindent\textit{Editorial note.} Preamble: the retained preamble of this chapter (sources/stacks-cotangent/retained-preamble.tex, lines 1-63), containing the Xy-pic setup, the comment environments, the theorem declarations and the macros including \textbackslash{}NL.\par\medskip
\noindent\textit{Allowed standard prerequisites (editorial).}\begin{itemize}
\item The computation of the lower shriek as sheafification of an explicit presheaf complex. Remark remark-compute-L-pi-shriek-spaces, used at proof lines 4727-4740. Describes Lpi\_! of the sheaf of differentials.
\item The section-level map identifying the complex with the ring-level cotangent complex. Remark remark-map-sections-over-U, used at proof lines 4740-4744. Identifies the complex as computing the cotangent complex of Lambda -> O\_X.
\item The description of the cotangent complex of an algebraic space over a ring. Lemma lemma-space-over-ring, used at proof lines 4744-4745. Concludes the identification with L\_\{X/Lambda\}.
\item Flatness of the module of differentials for the objects of the category. used at proof lines 4718-4724. Gives the second and third equalities of the displayed chain.
\end{itemize}\medskip
\ifdefined\hypertarget\hypertarget{sourcepointer-1}{}\fi
\noindent\textit{Source reference.} Original lemma 08VG: focused-seeds/AGP/sources/stacks-snapshot-20261011/tags.txt, source lines 11234--11234. Navigation pointer into the Stacks Project; the standard fact it supplies is applied inside the counted proof span and the argument resting on it stays there.\par\medskip
\ifdefined\hypertarget\hypertarget{sourcepointer-2}{}\fi
\noindent\textit{Source reference.} Original remark 08VL: focused-seeds/AGP/sources/stacks-snapshot-20261011/tags.txt, source lines 11239--11239. Navigation pointer into the Stacks Project; the standard fact it supplies is applied inside the counted proof span and the argument resting on it stays there.\par\medskip
\ifdefined\hypertarget\hypertarget{sourcepointer-3}{}\fi
\noindent\textit{Source reference.} Original remark 08SX: focused-seeds/AGP/sources/stacks-snapshot-20261011/tags.txt, source lines 11145--11145. Navigation pointer into the Stacks Project; the standard fact it supplies is applied inside the counted proof span and the argument resting on it stays there.\par\medskip
\begin{lemma}
\label{lemma-cotangent-morphism-spaces}
In the situation above there is a canonical isomorphism
$$
L_{X/\Lambda} = 
L\pi_!(Li^*\Omega_{\mathcal{O}/\underline{\Lambda}}) =
L\pi_!(i^*\Omega_{\mathcal{O}/\underline{\Lambda}}) =
L\pi_!(\Omega_{\mathcal{O}/\underline{\Lambda}}
\otimes_\mathcal{O} \underline{\mathcal{O}}_X)
$$
in $D(\mathcal{O}_X)$.
\end{lemma}

\begin{proof}
We first observe that for any object $(U \to \mathbf{A})$ of
$\mathcal{C}_{X/\Lambda}$
the value of the sheaf $\mathcal{O}$ is a polynomial algebra over $\Lambda$.
Hence $\Omega_{\mathcal{O}/\underline{\Lambda}}$ is a flat $\mathcal{O}$-module
and we conclude the second and third equalities of the statement of the
lemma hold.

\medskip\noindent
By Remark \ref{remark-compute-L-pi-shriek-spaces} the object
$L\pi_!(\Omega_{\mathcal{O}/\underline{\Lambda}}
\otimes_\mathcal{O} \underline{\mathcal{O}}_X)$
is computed as the sheafification of the complex of presheaves
$$
U \mapsto
\left(\Omega_{\mathcal{O}/\underline{\Lambda}}
\otimes_\mathcal{O} \underline{\mathcal{O}}_X\right)(\mathbf{A}_{\bullet, U})
=
\Omega_{P_{\bullet, U}/\Lambda} \otimes_{P_{\bullet, U}} \mathcal{O}_X(U) =
L_{\mathcal{O}_X(U)/\Lambda}
$$
using notation as in Remark \ref{remark-compute-L-pi-shriek-spaces}.
Now Remark \ref{remark-map-sections-over-U} shows that
$L\pi_!(\Omega_{\mathcal{O}/\underline{\Lambda}}
\otimes_\mathcal{O} \underline{\mathcal{O}}_X)$
computes the cotangent complex of the map of rings
$\underline{\Lambda} \to \mathcal{O}_X$ on $X_\etale$.
This is what we want by Lemma \ref{lemma-space-over-ring}.
\end{proof}

\end{document}
